\problemname{The Generous Traveler}
Your rich friend Erika lives in a country with $N$ villages where the $i$:th village has $A_i$ inhabitants.
The villages are connected with $N-1$ roads such that it is possible to travel from any village to any other
village using these roads.

After a whole day of work, Erika is exhausted and needs to go on vacation, so she plans to visit some of
the country's villages. As with all rich people, there is nothing your friend likes more than giving away money!
Therefore, she plans to take a big bag of coins and give away as much money as she can in each village
she visits (including the one she starts in).

However, there is a problem: if an inhabitant of a village receives less money than another inhabitant in
the same village, they become very angry, and Erika's life is threatened. To avoid this, she decides to
give all inhabitants as much money as she can so that everyone $\textit{in the same village}$ gets the
same amount, even if it means no one gets any money.

Now, Erika wonders how much money she will have left at the end of such a trip. She asks you $Q$ questions
of the form: Given that the trip starts at village $u$ with $x$ coins, and ends at village $v$, how much
money will be left at the end of the trip?

Note that the trip always follows the simple path between two villages, i.e., the unique shortest path.
Erika only gives out whole coins, so she always gives a non-negative integer amount of money to each villager.

\section*{Input}
The first line contains the two integers $N$ ($1 \leq N \leq 2 \cdot 10^5$) and $Q$ ($1 \leq Q \leq 10^5$).

Then follow $N-1$ lines, where the $i$-th line contains two integers $a_i$ and $b_i$ ($1 \leq a_i, b_i \leq N$).
This means that there is a road between villages $a_i$ and $b_i$.

The next line contains $N$ integers $A_1, A_2, \dots, A_N$ ($1 \leq A_i \leq 10^9$), each describing the number
of inhabitants in village $i$.

Finally, there are $Q$ lines, one per query.
The $i$-th line consists of integers $u_i$, $v_i$, $x_i$ ($1 \leq u_i, v_i \leq N$, $1 \leq x_i \leq 10^9$),
where $u_i$ is the starting village, $v_i$ is the ending village, and $x_i$ is the number of coins Erika
starts with in the $i$-th query.

\section*{Output}
For each query, print the answer on its own line.

\section*{Scoring}
Your solution will be tested on a set of test groups, each worth a number of points. Each test group contains
a set of test cases. To get the points for a test group you need to solve all test cases in the test group.

\noindent
\begin{tabular}{| l | l | p{12cm} |}
  \hline
  \textbf{Group} & \textbf{Points} & \textbf{Constraints} \\ \hline
  $1$   & $10$       & $N,Q \leq 2000$\\ \hline
  $2$   & $20$       & There is a path between village $i$ and village $i+1$ for all $1 \leq i \leq N-1$, $A_i \leq 100$. \\ \hline
  $3$   & $25$       & There is a path between village $i$ and village $i+1$ for all $1 \leq i \leq N-1$. \\ \hline
  $4$   & $25$       & All trips end at village 1. Formally, $v_i=1$ for all $1 \leq i \leq Q$. \\ \hline
  $5$   & $20$       & No additional constraints. \\ \hline
\end{tabular}


\section*{Explanation of sample}
In the first test case, the first journey occurs between village $8$ and village $10$, and Erika starts with $31$ kronor.
Already in the first village (with $20$ inhabitants), she gives away $1$ krona to each person, and has $11$ left.
The journey then continues to village $7$, with $10$ inhabitants. Here too, Erika gives away one krona to each person, and now only has one krona left.
Finally, Erika reaches village $10$, with $20$ inhabitants. Unfortunately, Erika can't afford to give any money to the population here, so she ends up with $1$ krona left.

The second test case is an example of a test case that can appear in test case groups 1, 2, and 3.

